\texttt{double.c} is a minimal C program used to generate the example LLVManim output. It defines three functions---\texttt{init}, \texttt{double\_1}, and \texttt{main}---whose control flow, heap allocation, and loop structure exercise all four LLVManim animation modes.
\inputminted[
  frame=single,
  bgcolor={rgb}{0.95,0.97,1.0}, 
  breaklines=true,
  fontsize=\footnotesize,
  linenos,
  xleftmargin=2em,
  numbersep=10pt,
]{c}{code_examples/double.c}
\captionof{listing}{Full C source (\texttt{double.c}) used for all sample animations.}
\label{lst:double.c}

\vspace{2em}
\needspace{6\baselineskip}

\texttt{double.ll} is compiled at \texttt{-O0} with full debug info (\texttt{clang -O0 -g -S -emit-llvm double.c}). It is the actual LLVM~IR input used by LLVManim in the examples.
\inputminted[
  frame=single,
  bgcolor={rgb}{0.95,0.97,1.0},
  breaklines=true,
  fontsize=\footnotesize,
  linenos,
  xleftmargin=2em, % Keep line numbers inside margins
  numbersep=10pt,
]{llvm}{code_examples/double.ll}
\captionof{listing}{Full LLVM~IR source (\texttt{double.ll}) used for all sample animations.}
\label{lst:double.ll}
